% Exact theorem excerpts, with mathematical notation transcribed into TeX.
% Retrieved 2026-08-24 from the publisher-hosted Numdam PDFs.

% A. J. de Jong, "Crystalline Dieudonne module theory via formal and
% rigid geometry", Publications Mathematiques de l'IHES 82 (1995), 5--96.
% DOI: 10.1007/BF02698637
% Source: https://www.numdam.org/item/10.1007/BF02698637.pdf

% Section 7 setup, pp. 74--75.
Let $O$ be a complete discrete valuation ring, $K$ its quotient field,
$\pi\in O$ a uniformizer and $k$ the residue field: $k=O/\pi O$.

We write $FS_O$ for the category of locally Noetherian adic formal schemes
$X$ over $\operatorname{Spf}(O)$ whose reduction $X_{\mathrm{red}}$ is a
scheme locally of finite type over $\operatorname{Spec}(k)$.

Let $X=\operatorname{Spf}(A)$ be an affine formal scheme over
$\operatorname{Spf}(O)$. Let $I\subset A$ be the biggest ideal of definition
of $A$. The topological ring $A$ is Noetherian and adic. Thus $X$ is an
object of $FS_O$ if and only if $A/I$ is a finitely generated $k$-algebra.
This is also equivalent to the fact that $A$ is a quotient of
\[
 O\{x_1,\ldots,x_n\}[[y_1,\ldots,y_m]]
\]
for some $n,m\in\mathbb N$.

% Lemma 7.1.9, pp. 78--79.
\textbf{7.1.9. Lemma.} Let $X=\operatorname{Spf}(A)$ be as above. There is a
bijection functorial in $A$ between the following two sets:
\begin{enumerate}
\item Maximal ideals $\mathfrak m\subset A\otimes_O K$.
\item Points $x$ of $X^{\mathrm{rig}}$.
\end{enumerate}
If the point $x(\mathfrak m)\in X^{\mathrm{rig}}$ corresponds to
$\mathfrak m\subset A\otimes_O K$, then there exists a canonical
homomorphism of local rings
\[
 (A\otimes_O K)_{\mathfrak m}\longrightarrow
 \mathcal O_{X^{\mathrm{rig}},x}
\]
compatible with the canonical map from $A\otimes_O K$ to global analytic
functions. This homomorphism induces an isomorphism on completions.

% The residue-field sentence in the proof of Lemma 7.1.9, p. 78.
If $x$ is represented by a maximal ideal $\mathfrak p$ in one of the
affinoid rings $C_n$, then the map $A\otimes_O K\to C_n/\mathfrak p$ is
surjective. Thus $\mathfrak m$ is its kernel and has the same residue field
as $\mathfrak p$ in $C_n$.

% Brian Conrad, "Irreducible components of rigid spaces",
% Annales de l'Institut Fourier 49 (1999), no. 2, 473--541.
% DOI: 10.5802/aif.1681
% Source: https://www.numdam.org/item/10.5802/aif.1681.pdf

% Formal-scheme convention, p. 476.
When the valuation on the complete field $k$ is discrete, with valuation
ring $R$, the class $FS_R$ consists of formal schemes over
$\operatorname{Spf}(R)$ which are locally noetherian adic and for which the
closed formal subscheme cut out by topologically nilpotent functions is
locally of finite type over the residue field of $R$.

% Theorem 2.1.3, p. 491.
\textbf{Theorem 2.1.3.} If $X$ is a locally finite type scheme over
$\operatorname{Spec}(k)$ and $f:\widetilde X\to X$ is a normalization, then
$f^{\mathrm{an}}$ is a normalization of $X^{\mathrm{an}}$. Likewise, if $k$
has discrete valuation and $X$ is a formal scheme in $FS_R$, and
$f:\widetilde X\to X$ is a normalization, then $f^{\mathrm{rig}}$ is a
normalization of $X^{\mathrm{rig}}$.

% Definition 2.2.2, p. 496.
\textbf{Definition 2.2.2.} The irreducible components of a rigid space $X$
are the reduced analytic sets of the form $X_i=p(\widetilde X_i)$, with
$\widetilde X_i$ the connected components of the normalization
$p:\widetilde X\to X$.
